% BEGIN THERMAL_R03_SYNC_01
\begingroup
\setlength{\parskip}{1.2pt plus .5pt minus .3pt}
\section{Introduction}

\subsection{Determination of constants and relations between action, area, and information}

This article determines the Planck action sections and their return,
the Barbero--Immirzi functional, and the CKM angular coordinates, and
establishes Boltzmann transduction and gravitational reciprocity in
the specified dimensional realizations. These results are connected
by effective operations: action converts frequency into energy;
Barbero--Immirzi weights the incidence that resolves area; Boltzmann
converts dimensionless information into entropy; and gravitation
transforms action into a common geometric scale. The physical question
is how these functions and their values are determined from a single
discrete structure instead of entering as measured parameters.

The determination of Planck separates the decimal order of action from
its return coordinates. The local quotient of the record has sixteen
sheets; its multiplicative norm, followed by autoscaling, gives
\(\operatorname{ord}_{10}(\varphi\alpha^{16})=-34\).
Regional characters then fix
\(\hbar_{\rm pre}=H_5\,10^{-34}\mathcal U_S\) and
\(\hbar_{\rm ret}=\eta_{\rm ret}\,10^{-34}\mathcal U_S\),
with \(0<\eta_{\rm ret}<H_5\). The ratio
\(R_{\rm act}=H_5/\eta_{\rm ret}>1\) preserves the effect of
return under a change of the common positive basis; in each section,
\(h_a=2\pi\hbar_a\). The local reader, the character, and the regional
continuation rule are part of this determination, developed
in \ref{sec:action}.

The Barbero--Immirzi functional evaluates an oriented chain of twelve
sectors and a return seam on incidences of cardinalities 90
and 120. Evaluation yields
\(\gamma_{\HMT}=0.2740076726944592747\ldots\), odd under angular
reversal. The CKM sector rules, in turn, produce a rational matrix
of determinant \(-3857/6\), three mixing angles, and the phase
\(\delta_{\rm deg}=9A_{\rm deg}\).
Their complex realization is unitary and determines the Jarlskog
invariant. Both results receive the angular coordinates already
generated: the Barbero functional does not select the mixing angles,
nor do those angles retrospectively select the sector rules.

Longitudinal composition precedes gravitational evaluation.
Incidence inscription and the complementary archive yield
\(r_N=5759/23040\). Longitudinal rule R4 applies this coefficient
to \(\alpha_{\HMT}^{16}L_*\) and determines \(L\); metric rule
R5 determines the positive radial operator of the same character used
by the energy. Section~\ref{g26:sec:rigidez-radial-en} states both
compositions and proves uniqueness of their joint realization with
phase action, the conjugate clock, and reduced-radius identification R8.
The reference \(L_*\), declared equal to one metre in the SI
evaluation, retains its dimensional role.

The relation between Boltzmann and gravitation then takes the explicit form
\[
 k_B\Theta_{{\rm clk},a}\log3=\frac{h_a}{T_{108}},\qquad
 G_a=\frac{c_{\rm int}^3R_{108}^2}{\hbar_a},\qquad
 \frac{G_a\hbar_a}{c_{\rm int}^3}=\ell_P^2.
\]
Here \(R_{108}=L\) and \(t_0=\pi_{\HMT}L/(54c_{\rm int})\).
The first identity determines the energy per tritic decision and a
unique thermal transducer once the positive temperature section is
fixed. The second is the circular specialization of the common
coefficient determined by the radial proof. The third shows that reciprocity
between action and gravitation preserves the area unit. These thermal
and geometric conditions relate the results without identifying
determination of a dimensional section with the choice of its numerical
coordinate.

The distinction requires attention to the nature of each object. A
dimensionless ratio expresses a relation between quantities. A dimensional
constant occupies a line of quantities and acquires a numerical coordinate
when a unit is chosen. A boundary condition fixes a realization within
a collection of solutions. All three determinations may enter
the same formula and preserve different roles in its proof.
Duff, Okun, and Veneziano's debate on dimensional constants
illustrates the importance of making these choices explicit
\cite{duffokunveneziano}. The current formulation of the International System
in turn specifies the role of defining constants in the
realization of units \cite{bipm2018,bipmbrochure}.

\subsection{Arithmetic support, orientation, and generating dynamics}

In Triadic Modular Holography, HMT, arithmetic support, ternary
orientation, and dynamics with memory precede scalar readings.
APP, \emph{All Possible Paths}, uses the discrete support
\(X_9=(\mathbb Z/9\mathbb Z)^2\), with the additive Cayley graph
generated by \(\pm(1,0)\) and \(\pm(0,1)\). At each position
of its \(1,\ldots,9\) chart, the two evaluations, sum and
product, act while retaining their residues and quotients. TRIT assigns
a signature \(\{-1,0,+1\}\) distinguishing orientation and local regime.
TPK, \emph{Topological Phase Kernel}, composes selection, transport,
carry and memory updates, and compatible state prolongation.
Phase return preserves memory advance. The operations and regional
generation in \ref{sec:nucleo}--\ref{sec:eta} construct the numerical
regions, the dodecaphase record, and alpha before their use in
action, incidence, and information readers.
The correlated analytic representation of alpha prolongs, at every
depth, the coordinate produced by dodecaphase closure.
The Archimedean realization determines values of compatible cylinders;
the incidence realization preserves other information from the same origin.
This common dependency allows each constant to be studied as value and
structure, with measurement subsequently serving as a comparison.

This distinction already takes material form in the constitution of the record.
The evaluation of visits and oriented traces in
\ref{subsec:registro-evaluacion-incidencial} produces three integer counts
per window, preserves their residues and quotients, and determines the twenty-five
coordinates of \eqref{eq:registro-composicion-incidencial}. The domain
includes the incidence collection, its multiplicities, and its orientation.
Proposition~\ref{prop:registro-36-residuos} characterizes exactly the
residual information determining that image; the transformations
in \ref{subsec:k-pi-h} and \ref{subsec:k-hadamard} subsequently allow the
signed record to be recovered. Evaluation of the ledger and
reconstruction of its output are thus linked, with a distinct demonstrative role
for each operation.

This structure acts effectively in the equations.
Incidence cardinalities enter a
return functional; action sections enter a
gravitational reciprocity; boundary occupations enter area
and the modular spectrum; inversion of a quadratic form enters
the exchange of momentum and winding. Each connection is expressed
through a map, an identity, and a domain of validity.

\subsection{Action, incidence, and boundary information}

The specific argument of Article II focuses on the relation between
action, area, entropy, and gravity. The Barbero--Immirzi parameter occupies
a central position because it links the incidence functional to an
areal normalization and its realization in gravitational action.
Holst's formulation provides a language of subsequent realization
for this coefficient \cite{holst1996}. In the construction studied here,
the coefficient is evaluated from oriented incidence and the angular
channels already obtained.

The inclusion of a marked hexad in an octad preserves the incidence
pair determining cardinalities 90 and 120. The dodecaphase turn
distinguishes twelve sectors from an oriented return boundary. The linear
reader of that chain produces the functional's multiplicities. The
pole coefficient and decimal evaluation are computed afterward.
The sequence reveals why the Barbero--Immirzi value belongs to a
structure of orientation and area, in addition to occupying a position in
a table of constants.

Area and the modular Hamiltonian describe complementary aspects
of the boundary. The first reading aggregates occupations. The second weights
their sectorial provenance. For \(\Delta_{\rm mod}\ne0\) and
known normalization, the Hamiltonian yields the weighted observable
\(D\); together with the total \(N\), it reconstructs the two
occupations originating the aggregate. Inversion of the pair
\((N,D)\) does not require that additional condition. This recoverability specifies
the content of the genealogical thesis: two configurations can
share a scalar reading and preserve differences accessible to another
reading of the same state. The argument is developed first in the
corresponding finite space and then in the declared realization.

Memory preservation also specifies how to reuse dynamics
when auxiliary variables are eliminated. Identity
\eqref{mem:identidad-finita} retains the terminal state of each truncation;
Proposition~\ref{mem:prop-serie} provides the prolongation and its tail
bound. Proposition~\ref{mem:prop-schur} proves that the complete
observation is preserved by jointly transforming the reduced operator,
source, and reader. This rule identifies the operative content that
must be transmitted between two descriptions of the same transport.

Entanglement entropy additionally requires a bipartition and a
state. The exposition specifies both through a purification and its
reduced state. The spectrum's dimensionless information is distinguished
from dimensional entropy, whose normalization incorporates the Boltzmann
constant. This precision allows the energy of tritic
periodicity to be related to action without confusing a state count,
a probability distribution, and a temperature.

\subsection{Capacity, energy-bearing memory, and thermometric realization}

Section~\ref{sec:ii-ter27-antecedentes} develops the finite mechanism
connecting renewal of the two APP sheets, capacity and the first vacancy,
double grading of the carrier, and memory energy. The bilateral loop
determines a decimal spectrum; tritic grouping retains the remainder.
The capacity marker and incidence sectors then produce a complete
intensity with registered energy origin. Connected currents fix the
intersector gaps, and centered transport preserves the reader between
realizations with different scales.

Section~\ref{sec:ii-termica-marcada} composes these results with action,
area, Barbero, and horizon. It adds a thermometric realization retaining
the spectral reference in a tensor record: the entropy channel reads
marked information, and the energy channel reads the conditioned mean.
The proof obtains the coefficient from their differential ratio and
establishes the corresponding variational principle. The fixed reference
and marked--conditioned pair of functionals are the constitutive conditions
of the realization. The chronological capacity Hamiltonian retains degrees
$23,26,29$ and determines its normalized state; thermal application of the
intensity is thereby related to clock and horizon, preserving their scales
and effective transports.

\subsection{Longitudinal composition and determination of gravitational coupling}

Gravitation forms part of the central argument. The composition of
inscription, length, and metric first fixes \(L\) and the radial response.
The conjugate clock satisfies \(\omega_{108}L=c_{\rm int}\); its period
is distributed over the 108 steps of the dodecaphasic calendar. This gives
\(T_{108}=2\pi_{\HMT}L/c_{\rm int}\) and \(t_0=T_{108}/108\).
Reduced-radius condition R8 determines \(G\) on the same character
already expressing the energy. Circular coincidence retains the equivalent
form \(\vartheta_{108}^{\circ}=(54/\pi_{\HMT})^2\). In the resulting
action sections, the gravitational coefficient transforms reciprocally,
and the product of action and gravitation remains invariant. The area
unit is \(L^2\) in both orientations.

The preserved equality has a geometric consequence:
passage between the sections affects certain coordinates while
preserving a common geometric quantity. The meaning of the value is
obtained from the transport that produces it. The relation to the horizon
then introduces energy, temperature, and areal entropy under
explicit normalizations. The confluence of these quantities in an
action cell is proved by substitution of the preceding
constructions.

Thermal dependency is presented with equal precision. An identity
for the product of the Boltzmann constant and a temperature
determines an energy scale. Determining each individual coordinate
requires its corresponding chart. The normalization of each realization specifies the bases and the
dimensional transport of that coordinate.
Physical comparison is performed on these outputs and their units.

\subsection{Reciprocity, angular propagation, and exceptional continuity}

Angular coordinates enter two distinct realizations
of the same construction. The oriented-incidence functional determines
the Barbero--Immirzi parameter; the species reader determines the
CKM mixing angles and their phase. Section~\ref{sec:ii-propagacion-angular}
presents the rational matrix, its inverse, and the complex unitary realization.
The determinant of the transformation and the Jarlskog invariant
provide complementary controls: the former examines recovery
of the coordinates; the latter examines the phase of the amplitudes. The
numerical values are compared after those constructions.

Theorem~\ref{thm:ii-witt-angular-transport} realizes the angular
coordinates in the Witt incidence module through an explicit
isometry. Corollary~\ref{cor:ii-witt-composed-reader} preserves the
metric of the sectorial system and provides an inverse on its image.
Thus passage to exceptional incidence transmits the coordinates
produced by the angular rules, together with their metric relations.
Sectorial evaluation and isometric transport compose in that
order, according to \eqref{eq:ii-witt-composed-reader}.

The action ellipse introduces a fixed-determinant form and an
anisotropy parameter. Exchanging its semiaxes inverts the
modular radius. Two transports can preserve the same energy
form and act with different signs on the symplectic form.
The orientation of action then distinguishes symplectic and
antisymplectic transports with the same scalar energy. This property underlies
the attention to transport and memory in the model's continuation.

The action--covariance relation specifies the quantum content of this
geometry. Proposition~\ref{prop:ii-covarianza-gaussiana} constructs,
in the canonical representation and with the Gaussian state specified there,
a covariance of determinant \(\hbar^2/4\) attaining equality in
Robertson--Schrödinger. Proposition~\ref{prop:ii-area-covarianza-reciprocidad}
identifies the action contour with the two-standard-deviation
ellipse: its semiaxes are \(2\Delta Q\) and \(2\Delta P\), and its area
is \(h=2\pi\hbar\). Reciprocity exchanges dispersions and
preserves that determinant. Passage to position and momentum keeps the
dimensional chart explicit, according to~\eqref{eq:ii-covarianza-dimensiones}.
The previously constructed action, its angular distribution,
and the fluctuations of a defined operator realization are thereby related.

Section~\ref{delta22:ii:transporte} proves exact conservation of
quadratic action under transport between ellipses, derives its
Hamiltonian connection term, and establishes the positivity condition
of the generator. The conformally symplectic version preserves
normalized action and specifies the transported form.


Momentum--winding exchange prolongs this reciprocity in
a lattice space with radius. Its proof preserves the carry
quotient, operator domain, and unitary transformation.
The corresponding modular function is evaluated on the radius thus
transported. Exceptional incidence is developed through to Moonshine
as part of the second realization of the common structure. The
connections between these realizations are presented through their concrete
maps, allowing what changes and what is preserved to be distinguished.

\Needspace{7\baselineskip}
Comparison of the quark and lepton sectors requires separately following
the antecedents that the angular development and species operators
transmit to CKM and PMNS. Each transport preserves its own frame; this
distinction makes it possible to specify where amplitude conjugation,
route orientation, and mixing invariants act. The manuscript
maintains construction differences in passing from structural
coordinates to the respective realizations.

\subsection{Identities for area, entropy, and gravitational scale}

The relation between information and gravitation is expressed through three
articulated determinations. The areal unit satisfies
$\ell_P^2=G\hbar/c^3$. Oriented incidence organizes its elementary
resolution through $\gamma_{\HMT}a_0$. The reduced state and bipartition
determine accessible information, whose dimensional expression uses
$k_B$. In the sectorial collection developed in
\ref{sec:ii-ecuacion-informacion}, these dependencies take the form
\begin{equation}
 \langle\widehat{\mathcal A}\rangle
 =\gamma_{\HMT}a_0\frac{G\hbar}{c^3}\langle N\rangle,
 \qquad
 S=k_B\bigl(\log Z+\Delta_{\rm mod}\langle D\rangle\bigr).
 \label{eq:ii-intro-area-informacion}
\end{equation}
The quantities $N$ and $D$ are the two readings of boundary
occupations; $Z$ normalizes the state. Their joint map recovers those
occupations through~\eqref{eq:recover}. The adjacent proofs also show
how their fluctuations govern the responses of area and
entropy.

The next step preserves its explicit geometric data. In the
horizon realization of Section~\ref{sec:confluencia},
$S_H/k_B=A_H/(4\ell_P^2)$ and $T_HS_H=E_\Lambda$.
Microscopic area, reduced-state entropy, and horizon
entropy retain their definitions and realization maps.
This organization allows incidence, action,
information, and gravitation to be read continuously, specifying at each transition the object
preserved and the observable that changes.

The scalar gravitational section also has a tensor realization.
Incidence descent and the rank-five projector in
\ref{sec:ii-pentadica} determine a five-component carrier;
its intertwiner realizes it as symmetric trace-free tensors before
the rank-two transverse reduction. Component orientation and
reduced observation thus remain distinct from the common coupling.
Microscopic area and sector information, constructed in
\ref{sec:area} and \ref{sec:ii-ecuacion-informacion}, retain their
definitions when these transports are composed. Exceptional continuity
in \ref{exc:doble-lectura} and \ref{exc:leech} preserves incidence
up to Moonshine; the duality in \ref{sec:ii-dualidad-t} preserves
the unitary action and operator domains under momentum--winding
exchange.

\par\endgroup
% END THERMAL_R03_SYNC_01
